%==============================================================================
%  James W. Kingston -- Curriculum Vitae
%  Reconstruction of KingstonCV.pdf in LaTeX.
%  Compiles with pdfLaTeX on Overleaf out of the box (no external files needed).
%==============================================================================
\documentclass[11pt,letterpaper]{article}

\usepackage[T1]{fontenc}
\usepackage[utf8]{inputenc}
\usepackage{mathptmx}                 % Times-like serif, matches the Word original
\usepackage[margin=1in]{geometry}
\usepackage{enumitem}
\usepackage{titlesec}
\usepackage{xcolor}
\usepackage[hidelinks]{hyperref}
\usepackage{needspace}

%------------------------------------------------------------------------------
% Section headings: large, grey, non-bold -- mimics the Word "Heading 1" style
%------------------------------------------------------------------------------
\definecolor{headinggray}{HTML}{595959}

\titleformat{\section}
  {\normalfont\Large\color{headinggray}}{}{0pt}{}
\titlespacing*{\section}{0pt}{1.2\baselineskip}{0.25\baselineskip}

%------------------------------------------------------------------------------
% List styles
%------------------------------------------------------------------------------
% First-level bullets use a dash, like the original.
\setlist[itemize,1]{label=--, leftmargin=2.2em, topsep=1pt, itemsep=1pt,
                    parsep=0pt, labelsep=0.6em}
\setlist[itemize,2]{label=\textrm{o}, leftmargin=2.0em, topsep=1pt, itemsep=1pt,
                    parsep=0pt, labelsep=0.6em}
\setlist[enumerate,1]{leftmargin=2.2em, topsep=2pt, itemsep=2pt, parsep=0pt}

%------------------------------------------------------------------------------
% Helper macros
%------------------------------------------------------------------------------
% \entry{Title}{Dates} -- bold title flush left, dates flush right
\newcommand{\entry}[2]{%
  \needspace{3\baselineskip}%
  \noindent\textbf{#1}\hfill #2\par}

% \affil{Institution, City ST ZIP} -- italic affiliation line under an entry
\newcommand{\affil}[1]{\noindent\textit{#1}\par\vspace{1pt}}

% \pubitem -- publication entries get a little breathing room
\newenvironment{publist}
  {\begin{enumerate}[leftmargin=2.2em, topsep=2pt, itemsep=6pt, parsep=0pt]}
  {\end{enumerate}}

\setlength{\parindent}{0pt}
\pagestyle{empty}

%==============================================================================
\begin{document}

%------------------------------------------------------------------------------
% Header
%------------------------------------------------------------------------------
\begin{center}
  {\Large James W. Kingston}\\[2pt]
  \href{mailto:james.kingston@yale.edu}{james.kingston@yale.edu}
\end{center}
\vspace{0.5\baselineskip}

%------------------------------------------------------------------------------
\section*{EDUCATION}

\entry{University of California, Davis}{2026}
\begin{itemize}
  \item Physics, PhD (GPA: 4.0)
\end{itemize}

\entry{University of Chicago}{2019}
\begin{itemize}
  \item Physics, Masters
\end{itemize}

\entry{University of California, Berkeley}{2017}
\begin{itemize}
  \item Physics and Applied Mathematics, Bachelors
\end{itemize}

%------------------------------------------------------------------------------
\section*{SKILLS}

\textbf{Operating Systems}: UNIX/Linux, Windows\\
\textbf{Programming Languages}: Python, C++, Bash\\
\textbf{Software}: NumPy, Matplotlib, SciPy, COMSOL, SolidWorks, ROOT, KiCad\\
\textbf{Hardware}: High Vacuum Systems, Noble Element Cryogenics, High Voltage (HV),
Electronics, Soldering, Machine Shop Proficiency, Orbital Welding

%------------------------------------------------------------------------------
\section*{PROFESSIONAL EXPERIENCE}

\entry{Graduate Student Researcher, Rare Event Detection Group}{May 2021 -- May 2026}
\affil{Lawrence Livermore National Laboratory, Livermore CA 94550}
\begin{itemize}
  \item Studied properties of electroluminescence light in the CHILLAX detector, the
        world's first dual phase xenon-doped argon time projection chamber (TPC),
        culminating in a first-author paper.
  \item Analyzed scintillation waveforms captured by Silicon Photomultipliers (SiPMs)
        from CHILLAX, including:
        \begin{itemize}
          \item Scrutinizing individual event waveforms to tune pulse-finding and
                baseline-finding algorithms.
          \item Analyzing reduced quantities (RQs) to perform data quality cuts.
          \item Performing statistical analyses on RQs to draw conclusions about data.
        \end{itemize}
  \item Designed and built a TPC for CHILLAX, which hosts a 35~g active target, 6 SiPMs,
        and over 6~kV/cm drift field.
  \item Set up analog electronic scheme for amplifying, filtering, and triggering SiPM
        signals allowing for single photoelectron (SPE) sensitivity in each channel.
  \item Developed and built cryogenics for CHILLAX which proved capable of stable
        operation at 5\% xenon concentration, resulting in a publication.
  \item Built a gas panel for purifying, circulating, precision doping, and mixing of
        noble gases.
  \item Designed and constructed an electronics rack for mounting and wiring all
        electronics for CHILLAX, which includes kempt routing of hundreds of wires for
        over 50 instruments.
  \item Designed and instrumented a cryogenic capacitor sensitive to 1~fF shifts
        corresponding to 0.01\% changes to the xenon concentration in liquid argon.
  \item Conceived and built a bubble routing device to direct undesired bubbles away
        from the detector liquid volume, stabilizing the liquid surface.
  \item Characterized SiPM dark rate, crosstalk, afterpulsing, and photon detection
        efficiency (PDE) in a liquid argon environment with direct argon scintillation
        light, resulting in a publication.
\end{itemize}

\vspace{0.4\baselineskip}
\entry{Assistant Specialist, Pantic Group}{September 2019 -- August 2020}
\affil{Department of Physics, UC Davis, Davis CA 95616}
\begin{itemize}
  \item Worked in a team to develop a mockup of the DarkSide-20k argon detector for HV
        studies.
  \item Performed a COMSOL electric field simulation to find the optimal HV cable
        properties for minimizing an electric field hotspot.
  \item Designed and built a 10$'$ crane frame to lift the mockup cryostat flange for
        maintenance.
  \item Designed a cryogenic camera for viewing HV sparking inside the cryostat.
\end{itemize}

\vspace{0.4\baselineskip}
\entry{Graduate Researcher, Grandi Group}{September 2018 -- June 2019}
\affil{Kavli Institute for Cosmological Physics, University of Chicago, Chicago IL 60615}
\begin{itemize}
  \item Operated and collected data with a dual-phase xenon TPC.
  \item Expanded on existing Geant4 simulations to measure geometric efficiency of
        photomultiplier tubes (PMTs).
  \item Designed and implemented a level meter inside the xenon detector for precision
        measurements of the gas-liquid interface.
  \item Analyzed data using techniques such as pulse-shape discrimination and event
        position-reconstruction.
  \item Calibrated the detector with krypton-83 gammas and measured the half-life of
        decays with 3\% precision.
\end{itemize}

\vspace{0.4\baselineskip}
\entry{Research Assistant, Rare Event Detection Group}{June 2016 -- August 2018}
\affil{Lawrence Livermore National Laboratory, Livermore CA 94550}
\begin{itemize}
  \item Worked in a team to design and build XeNu, a dual-phase xenon TPC.
  \item Made 3D potential and electric field simulations of the detector volume.
  \item Calibrated detector response to single photoelectrons and single electrons.
  \item Assisted in operating XeNu TPC HV at liquid drift fields up to 10.4~kV/cm for
        electron extraction efficiency measurements, the highest achieved to date in a
        dual-phase xenon detector.
  \item Commissioned detector at Triangle Universities Nuclear Laboratory to study low
        energy nuclear recoils from a neutron beam, where the ionization yield of
        0.3~keV nuclear recoils was measured for the first time.
\end{itemize}

\vspace{0.4\baselineskip}
\entry{Undergraduate Researcher, CMS Group}{July 2017 -- September 2017}
\affil{Deutsches Elektronen-Synchrotron (DESY), Hamburg Germany 22607}
\begin{itemize}
  \item Added commands to a program that sent instructions to and read data from a new
        pixel readout chip.
  \item Wrote code to send a calibration pulse to the readout chip and draw the
        resulting signal pulse shape.
  \item Characterized the chip by determining the dependence of features of the pulse
        shape (e.g.\ amplitude, peaking time) on various parameters of the chip that
        could be modified by the user.
  \item Optimized the chip for usage in the DESY electron test beam.
\end{itemize}

%------------------------------------------------------------------------------
\section*{CERTIFICATIONS}

\entry{Machine Shop Class}{April 2023 -- September 2023}
\affil{UC Davis Mathematical and Physical Science Division Machine Shop, Davis CA 95616}
\begin{itemize}
  \item Learned basic machine tool and shop safety, fundamentals of precision measuring,
        introduction to materials.
  \item Operated lathes, milling machines, and other powered shop equipment.
  \item Read and made computer-aided design (CAD) drawings.
  \item Applied skills to build a functional tape dispenser composed of self-machined
        parts made of aluminum, acrylic, and stainless steel.
\end{itemize}

%------------------------------------------------------------------------------
\section*{PUBLICATIONS}

\begin{publist}
  \item \textbf{J. Kingston}, J. Qi, J. Xu, E. Bernard, A. Tidball, A. Peck, N. Bowden,
        M. Tripathi, K. Ni, S. Westerdale, ``Gas electroluminescence in a dual phase
        xenon-doped argon detector'', Phys.\ Rev.\ D \textbf{113}, 032016
        (February 2026)

  \item J. Xu, D. Adams, B. Lenardo, T. Pershing, R. Mannino, E. Bernard,
        \textbf{J. Kingston}, E. Mizrachi, J. Lin, R. Essig, V. Mozin, P. Kerr,
        A. Bernstein, M. Tripathi, ``Search for the Migdal effect in liquid xenon with
        keV-level nuclear recoils'', Phys.\ Rev.\ D \textbf{109}, L051101 (March 2024)

  \item E. Bernard, E. Mizrachi, \textbf{J. Kingston}, J. Xu, S. Pereverzev,
        T. Pershing, R. Smith, C. Prior, N. S. Bowden, A. Bernstein, C. Hall, E. Pantic,
        M. Tripathi, D. McKinsey, P. Barbeau, ``Thermodynamic stability of xenon-doped
        liquid argon detectors'', Phys.\ Rev.\ C \textbf{108}, 045503 (October 2023)

  \item T. Pershing, D. Naim, B. Lenardo, J. Xu, \textbf{J. Kingston}, E. Mizrachi,
        V. Mozin, P. Kerr, S. Pereverzev, A. Bernstein, M. Tripathi, ``Calibrating the
        scintillation and ionization responses of xenon recoils for high-energy dark
        matter searches'', Phys.\ Rev.\ D \textbf{106}, 052013 (September 2022)

  \item T. Pershing, J. Xu, E. Bernard, \textbf{J. Kingston}, E. Mizrachi, J. Brodsky,
        A. Razeto, P. Kachru, A. Bernstein, E. Pantic, I. Jovanovic, ``Performance of
        Hamamatsu VUV4 SiPMs for detecting liquid argon scintillation'', Journal of
        Instrumentation \textbf{17}, P04017 (May 2022)

  \item B. Lenardo, J. Xu, S. Pereverzev, O. Akindele, D. Naim, \textbf{J. Kingston},
        A. Bernstein, K. Kazkaz, M. Tripathi, C. Awe, L. Li, J. Runge, S. Hedges, P. An,
        P. Barbeau, ``Low-energy physics reach of xenon detectors for
        nuclear-recoil-based dark matter and neutrino experiments'',
        Phys.\ Rev.\ Lett.\ \textbf{123}, 231106 (December 2019)

  \item J. Xu, S. Pereverzev, B. Lenardo, \textbf{J. Kingston}, D. Naim, A. Bernstein,
        K. Kazkaz, M. Tripathi, ``Electron extraction efficiency study for dual-phase
        xenon dark matter experiments'', Phys.\ Rev.\ D \textbf{99}, 103024 (May 2019)
\end{publist}

%------------------------------------------------------------------------------
\section*{ORAL PRESENTATIONS}

\begin{publist}
  \item ``Gas Electroluminescence in a Dual Phase Xenon-Doped Argon TPC'', Light
        Detection in Noble Elements (Hong Kong, October 2025)

  \item ``Recent Results of a Xenon-Doped Argon Ionization Detector with the CHILLAX
        Experiment'', American Physical Society (APS) Global Physics Summit
        (Anaheim, March 2025)

  \item ``Xenon-Doping of Liquid Argon with the CHILLAX Detector'', Nuclear Science and
        Security Fall Workshop (Livermore, November 2024)

  \item ``Reactor CEvNS: LXe and LAr Detectors'', Workshop on Neutrino Science and
        Applications at the High Flux Isotope Reactor (Oak Ridge, April 2024)

  \item ``Studies of Xenon-Doped Argon with the CHILLAX Experiment'', Coordinating Panel
        on Advanced Detectors (Menlo Park, November 2023)

  \item ``Capacitive Monitoring of Xenon Concentration in a Xenon-Doped Argon
        Detector'', Coordinating Panel on Advanced Detectors (Stony Brook,
        November 2022)

  \item ``Characterization of Hamamatsu VUV4 Silicon Photomultipliers using Argon
        Scintillation Light'', APS Grad Slam Lightning Talk \textbf{(Winner, Awarded
        1 Year of APS Free Membership)} (New York City, April 2022)

  \item ``Characterization of Hamamatsu VUV4 Silicon Photomultipliers using Argon
        Scintillation Light'', APS April Meeting (New York City, April 2022)
\end{publist}

%------------------------------------------------------------------------------
\section*{FELLOWSHIPS AND AWARDS}

\begin{enumerate}
  \item Professor Emerita Ling-Lie Chau Prize for Excellence in Graduate Studies
        (June 2025)
  \item Nuclear Science and Security Consortium (NSSC) Best National Laboratory
        Collaboration Award (June 2024)
  \item UC Davis Graduate Program Fellowship (Spring 2024)
  \item High Energy Physics Consortium for Advanced Training (HEPCAT) Fellowship
        (2023--2024)
  \item NSSC Graduate Fellow (2021--2026)
  \item NSSC Undergraduate Fellow (2016--2017)
\end{enumerate}

%------------------------------------------------------------------------------
\section*{OUTREACH}

\begin{enumerate}
  \item Lab Instructor, HEPCAT Summer School \hfill Summer 2022, 2023, 2024
        \begin{itemize}
          \item Deployed a minimal liquid argon scintillation detector for demonstration
                at the HEPCAT Summer School.
          \item Instructed students on basics of safely operating vacuum, high pressure,
                and cryogenic hardware.
          \item Guided multiple small groups of students to submerge the detector's
                SiPMs in approximately half a liter of liquid argon.
          \item Utilized a radioactive source to increase the detector event rate and
                demonstrate the SiPMs' sensitivity to liquid argon scintillation light.
        \end{itemize}
\end{enumerate}

%------------------------------------------------------------------------------
\section*{TEACHING EXPERIENCE}

\begin{enumerate}
  \item Physics 80, UC Davis \hfill Spring 2021
        \begin{itemize}
          \item Taught a lab section for a Python data analysis course for physics
                majors.
          \item Wrote solutions for and graded assignments.
        \end{itemize}

  \item Physics 9B, UC Davis \hfill Winter 2021
        \begin{itemize}
          \item Held weekly office hours for an introductory electromagnetism course for
                physics majors.
          \item Wrote solutions for and graded assignments.
        \end{itemize}

  \item Physics 7A, UC Davis \hfill Fall 2020
        \begin{itemize}
          \item Taught a lab and discussion section for an introductory conservation
                principles course for non-physics majors.
          \item Graded lab assignments.
        \end{itemize}
\end{enumerate}

\end{document}
